\subsubsection{Applications of Autoencoders}\label{sec:applications-of-autoencoders}

Although at first glance an autoencoder (\autopageref{sec:neural-autoencoders}) may seem useless (after all, it is trained to reproduce its input) forcing the network to pass information through a compressed or constrained latent representation allows it to \textbf{learn useful structures hidden in the data}. Because of this property, autoencoders have many practical applications.
\begin{itemize}
    \item \important{Data Compression}. One of the most natural applications of autoencoders is \textbf{data compression}. The encoder transforms the original input vector $x \in \mathbb{R}^{I}$ into a lower-dimensional latent representation $h \in \mathbb{R}^{J}$, with $J < I$, which contains only the most relevant information needed to reconstruct the input. The decoder can later use this compressed representation to approximately recover the original data.

    Unlike traditional compression algorithms, the encoder and decoder are learned directly from data and can therefore capture complex nonlinear relationships.
    
    
    \item \important{Dimensionality Reduction}. Autoencoders can also be viewed as a form of \textbf{nonlinear dimensionality reduction}. Classical techniques such as Principal Component Analysis (PCA) project data onto a lower-dimensional linear subspace. Autoencoders generalize this idea by learning nonlinear mappings, allowing them to discover more complex structures in the data. For this reason, they are often described as a kind of nonlinear PCA.

    The latent representation learned by the encoder can then be used for visualization, clustering, or as input to other machine learning algorithms.
    
    
    \item \important{Feature Extraction}. The hidden representation learned by an \textbf{autoencoder often contains more meaningful information than the raw input}. Instead of using the original data directly, we can use the latent vector $h$ as a new set of features. These features typically capture higher-level patterns and correlations present in the data.
    
    Historically, autoencoders were frequently used as a pretraining mechanism: a network was first trained in an unsupervised way to learn useful features and then fine-tuned on a supervised task.
    
    
    \item \important{Denoising}. Autoencoders can be trained to \textbf{remove noise from data}. Instead of reconstructing a clean input from itself, the network receives a corrupted version $\tilde{x}$ and is trained to reconstruct the original clean sample $x$. This forces the encoder to learn robust features that capture the underlying structure of the data rather than the noise itself. Such \textbf{denoising autoencoders} are widely used in image restoration, audio enhancement, and signal processing.
    
    
    \item \important{Anomaly Detection}. An autoencoder \textbf{trained only on normal examples becomes highly specialized at reconstructing them}.
    
    When an unusual or anomalous sample is presented, the network is typically unable to reconstruct it accurately because it does not match the patterns seen during training. Consequently, anomalous inputs tend to produce a much larger reconstruction error:
    \begin{equation*}
        \text{Reconstruction Error} = \| x - \hat{x} \|^2
    \end{equation*}
    Where $x$ is the original input and $\hat{x}$ is the reconstructed output. This property makes autoencoders useful for fraud detection, industrial monitoring, cybersecurity, and fault diagnosis.
    
    
    \item \important{Generative Models}. More advanced variants of autoencoders can be used to \textbf{generate new data}. The most famous example is the Variational Autoencoder (VAE), which learns a structured latent space from which new samples can be generated by drawing random latent vectors and decoding them. VAEs represent an important family of generative models and are the foundation of several modern generative techniques.
    
    
    \item \important{Word Embeddings} (our main focus in this section). Autoencoders have also \textbf{inspired methods for learning word embeddings}, which are \textbf{dense vector representations of words}.

    The main idea is conceptually similar: instead of representing a word as a very large sparse one-hot vector (i.e., a vector with a single 1 and the rest 0s, a word gets its own unique position in the embedding space), a neural network learns a compact latent representation in a much lower-dimensional continuous space. \hl{Similar words tend to be mapped to nearby points in this space.}
    
    For example, words such as \emph{king} and \emph{queen}, or \emph{Paris} and \emph{Rome}, often exhibit meaningful geometric relationships in the embedding space. Word embeddings will be studied in detail in the next chapter, where we will examine techniques such as Word2Vec and GloVe.
\end{itemize}